\chapter{배치, 샘플링, 검증: 데이터 흐름이 결과를 결정한다}
\label{bg:chapter-batch-evaluation}

\section{Batch는 단지 GPU를 채우는 단위가 아니다}

\concept{batch}{배치}는 한 update의 기울기를 추정하는 예시 집합이다. 전체 데이터의
평균 기울기를 매번 계산하는 대신 작은 subset의 평균을 사용하면 계산은 줄지만 추정에
noise가 생긴다. 적당한 noise는 평평한 해를 찾는 데 도움이 될 수 있고, 지나친 noise는
학습을 불안정하게 만든다. batch를 키우면 noise가 줄지만 activation memory와 통신량이
늘고, 같은 epoch에서 update 횟수가 줄어든다.

\concept{epoch}{에폭}은 데이터 전체를 대략 한 번 소비한 단위다. 고정 corpus에서는
직관적이지만, 계속 유입되는 agent experience에서는 명확한 epoch가 없다. 이때는 processed
tokens, unique episodes, wall-clock, user-days, consolidation cycles를 별도로 기록해야 한다.

\section{Sampling은 무엇을 기억하고 무엇을 잊을지 정한다}

\concept{sampling}{샘플링}은 다음 batch에 들어갈 데이터를 정한다. 균등 sampling은
흔한 사건을 반복 학습하고 rare but important event를 거의 보지 못할 수 있다. recency
sampling은 freshness를 높이지만 오래된 지식을 밀어낸다. surprise sampling은 모델이 크게
틀린 경험을 강조하지만 adversarial outlier와 noise에 취약하다. value sampling은 future
reuse가 높은 memory를 선택하려 하지만 그 value를 사전에 알기 어렵다.

Sleep system의 admission score를 개념적으로 다음처럼 둘 수 있다.

\begin{equation}
  s_i = \alpha u_i + \beta n_i + \gamma r_i
        -\delta q_i-\epsilon p_i,
  \label{eq:bg-admission-score}
\end{equation}

$u_i$는 관측된 utility, $n_i$는 novelty, $r_i$는 expected reuse, $q_i$는 품질 위험,
$p_i$는 privacy/governance 비용이다. 이것은 established formula가 아니라 scheduler가
숨겨 둔 판단을 명시적으로 드러내는 accounting form이다. 계수는 offline replay와
counterfactual evaluation으로 검증해야 한다.

\section{Training, validation, test의 역할}

\concept{validation-set}{검증 세트}는 hyperparameter, checkpoint, stopping 시점을 고르는 데
사용한다. \concept{test-set}{테스트 세트}는 선택을 마친 뒤 최종 일반화 성능을 추정한다.
test 결과를 보고 계속 모델을 수정하면 test도 사실상 validation이 되고, reported score는
낙관적으로 편향된다.

Memory system에서는 split 단위가 더 어렵다. 같은 사용자의 같은 사건을 paraphrase한 문장이
train과 test에 나뉘면 정보 누출이 생긴다. 공개 benchmark를 검색할 수 있는 agent는 test
question 자체를 sleep memory에 저장할 수 있다. 따라서 다음 split을 함께 고려해야 한다.

\begin{itemize}
  \item \textbf{time split}: 과거 episode로 학습하고 미래 interaction으로 평가한다.
  \item \textbf{entity split}: 특정 사람·프로젝트·도메인을 통째로 보존한다.
  \item \textbf{lineage split}: 동일 원문에서 파생된 summary와 paraphrase를 같은 split에 둔다.
  \item \textbf{policy split}: admission/eviction 정책의 학습 trace와 평가 trace를 분리한다.
  \item \textbf{contamination audit}: benchmark 문항이나 정답이 검색·기억에 들어왔는지 검사한다.
\end{itemize}

\section{Overfitting과 generalization}

\concept{overfitting}{과적합}은 training data에서 loss가 낮지만 새로운 data에서 성능이
나빠지는 현상이다. \concept{generalization}{일반화}는 보지 못한 상황에서도 배운 규칙이
유효한 정도다. 개인 memory에서는 ``같은 사실을 정확히 다시 말하기''와 ``그 사실을 새로운
작업에 적절히 적용하기''를 구분해야 한다. 전자는 retention, 후자는 transfer에 가깝다.

장기기억 평가는 최소 다섯 축이 필요하다.

\begin{enumerate}
  \item \textbf{write}: 올바른 정보가 destination state에 들어갔는가?
  \item \textbf{access}: 관련 query에서 그 정보가 실제 행동에 도달하는가?
  \item \textbf{retention}: 시간이 지나고 다른 update 뒤에도 유지되는가?
  \item \textbf{composition}: 다른 기억과 결합해 새로운 문제를 푸는가?
  \item \textbf{governance}: source를 추적하고 삭제·수정·rollback할 수 있는가?
\end{enumerate}

\begin{table}[htbp]
\centering
\caption{Sleep update의 평가 세트 예시}
\label{tab:bg-sleep-evaluation}
\begin{tabularx}{\textwidth}{L{0.20\textwidth}YYL{0.20\textwidth}}
\toprule
세트 & 포함 내용 & 실패가 뜻하는 것 & release gate \\
\midrule
new-memory query & 방금 통합한 사실·선호 & write/access 실패 & 핵심 recall 하한 \\
old-memory query & 이전 cycle의 기억 & catastrophic forgetting & regression 한도 \\
base capability & 범용 지식·추론·안전 & parametric side effect & non-inferiority \\
privacy/security & secret, poisoning, injection & 영구 오염 위험 & zero critical leak \\
systems & latency, bytes, energy, rollback & wake SLA/비용 실패 & budget ceiling \\
\bottomrule
\end{tabularx}
\end{table}

\section{통계적 유의성과 운영적 유의성}

표본이 크면 작은 차이도 통계적으로 유의할 수 있다. 하지만 0.1점 평균 향상이 adapter
수천 개, daily sleep cluster, 수십 TB write amplification을 정당화하지는 않는다. 반대로
희귀하지만 치명적인 deletion failure는 평균 score에 거의 보이지 않는다. 따라서 effect
size, confidence interval, tail risk, subgroup worst case, cost를 함께 보고한다.

\begin{keypoint}
Sleep-time compute의 올바른 test는 ``sleep 후 score가 올랐는가'' 하나가 아니다.
같은 future workload에서 external retrieval, longer context, cache, periodic refresh,
LoRA sleep을 동일한 data budget과 total cost로 비교해야 한다.
\end{keypoint}

